%
% 23-ableitungen.tex -- Ableitungen
%
% (c) 2026 Prof Dr Andreas Müller
%
\subsection{Ableitungen}
Der Satz~\ref{buch:integration:satz:residuensatz} zeigt, wie der
Wert einer holomorphen Funktion in einem Punkt $z_0$ allein aus den
Werten entlang eines Weges berechnet werden kann, der den Punkt $z_0$
im Gegenuhrzeigersinn umläuft.
Die durch die Formel
\eqref{buch:integration:eqn:residuensatz}
gegebene Funktion ist holomorph, man sollte daher auch die Ableitung
mit Hilfe eines Integrals berechnen können.
Wenn man 
\eqref{buch:integration:eqn:residuensatz}
unter dem Integral nach $z_0$ ableitet, erhält man
\begin{align}
f'(z_0)
=
\frac{d}{dz_0}f(z_0)
&=
\frac{1}{2\pi i}
\oint_{\gamma}
\frac{d}{dz_0}
\frac{f(z)}{z-z_0}
\,dz
\notag
\\
&=
\frac{1}{2\pi i}
\oint_{\gamma}
\frac{f(z)}{(z-z_0)^2}
\,dz.
\label{buch:integration:eqn:ableitung2}
\end{align}
Wie erwartet kann die Ableitung von $f$ ebenfalls durch ein einfaches
Wegintegral berechnet werden.
Durch wiederholtes Ableiten kann man auch höhere Ableitungen finden,
wie der folgende Satz zeigt.

\begin{satz}[Cauchy-Formel für Ableitungen]
\label{buch:integration:satz:ableitungn}
\index{Satz!Cauchy-Formel für Ableitungen}%
\index{Cauchy-Ableitungsformel}%
Ist $f\colon U\to\mathbb{C}$ eine holomorphe Funktion und $\gamma$
ein Weg, der den Punkt $z_0\in U$ im Gegenuhrzeigersinn umläuft.
Dann gilt
\begin{equation}
\frac{d^n}{dz_0^n}f(z_0)
=
f^{(n)}(z_0)
=
\frac{n!}{2\pi i}
\oint_{\gamma}
\frac{f(z)}{(z-z_0)^{n+1}}
\,dz.
\label{buch:integration:eqn:ableitungn}
\end{equation}
\end{satz}

\begin{proof}
Der Fall $n=1$ ist bereits durch \eqref{buch:integration:eqn:ableitung2}
bewiesen, wir müssen die Formel
\eqref{buch:integration:eqn:ableitungn}
für $n>1$ beweisen.
Wir tun dies mit Hilfe vollständiger Induktion, wobei 
\eqref{buch:integration:eqn:ableitung2} als Induktionsverankerung
dient.

Wir nehmen jetzt an \eqref{buch:integration:eqn:ableitungn} für $n$ gilt,
und zeigen, dass sie dann auch für $n+1$ gilt.
Dazu leiten wir \eqref{buch:integration:eqn:ableitungn} nach $z_0$
ab und erhalten
\begin{align*}
f^{(n+1)}(z_0)
&=
\frac{d}{dz_0}
f^{(n)}(z_0)
=
\frac{d}{dz_0}
\frac{n!}{2\pi i}
\oint
\frac{f(z)}{(z-z_0)^{n+1}}
\,dz
\\
&=
\frac{n!}{2\pi i}
\oint
\frac{d}{dz_0}
\frac{f(z)}{(z-z_0)^{n+1}}
\,dz
\\
&=
\frac{n!}{2\pi i}
\oint
(-(n+1))
\frac{f(z)}{(z-z_0)^{n+2}}
(-1)
\,dz
\\
&=
\frac{(n+1)!}{2\pi i}
\oint
\frac{f(z)}{(z-z_0)^{n+2}}
\,dz.
\end{align*}
Damit ist \eqref{buch:integration:eqn:ableitungn} bewiesen.
\end{proof}

Die Werte und alle Ableitungen einer holomorphen Funktion in der Umgebung
eines Punktes können also durch Berechnung von Integralen der Form
\eqref{buch:integration:eqn:ableitungn}
über einen Weg, der die Umgebung umschliesst, berechnet werden.
Da der Integrand entlang des Weges stetig ist, sind diese Integrale
alle wohldefiniert.
Da das Bild des Weges kompakt ist, sind die Integrale als Funktionen
von $z_0$ stetig.
Wir erhalten damit das folgende überraschende Resultat.

\begin{korollar}
Eine holomorphe Funktion $f\colon U\to\mathbb{C}$ ist beliebig oft stetig
differenzierbar.
\end{korollar}

Einmalige komplexe Differenzierbarkeit hat also automatisch beliebige
stetige Differenzierbarkeit zur Folge.
Dies steht in krassem Gegensatz zur Situation reeller Funktionen,
wo Beispiele von einmal stetig differenzierbaren Funktionen konstruiert
worden sind, die nicht beliebig oft differenzierbar sind.
Die komplexe Differenzierbarkeit ist also eine sehr viel stärkere
Einschränkung als reelle Differenzierbarkeit.

